\centerline{\bf \Large  РЕШЕНИЯ и КРИТЕРИИ}

\centerline{\bf \Large  11 КЛАСС}

\problem{\textbf{1}} \par

$$
\sqrt{2025+a \sqrt{2025+a \sqrt{2025+. .}}}=2025
$$


По условию задачи, когда количество радикалов стремится к бесконечности, мы можем записать:

$$
\sqrt{2025+a \sqrt{2025+a \sqrt{2025+. .}}}=2025
$$


Следовательно:

$$
\begin{gathered}
\sqrt{2025+a * 2025}=2025 \\
2025+a * 2025=2025^2 \\
a * 2025=2025(2025-1) \\
a=\frac{2025(2025-1)}{2025}=2024
\end{gathered}
$$


Ответ:

$$
a=2024
$$


\textbf{Критерии.} \\
\textit{7 баллов} $-$ Полностью верное решение\\
\textit{6 баллов} $-$ Правильное решение с арифметической ошибкой\\
\textit{0 баллов} $-$ Только ответ\\

\problem{\textbf{2}} Так как все треугольники равнобедренные и угол подъёма одинаков, это означает, что углы при основаниях этих равнобедренных треугольников равны. Значит, все равнобедренные треугольники подобны. Так как каждая последующая высота в 3 раза больше, то и сумма боковых сторон в 3 раза больше, соответственно, времени на переход через каждый последующий холм требуется в 3 раза больше. Подъём занимает 2 часа, а спуск - в два раза меньше. Значит, спускается за 1 час. То есть первый холм перейдет за 3 часа, второй - за $3\cdot3$ часа, третий - за $3\cdot3^2$ часов, ... , $n$-ый - за $3\cdot3^n$ часов. Вся эта сумма не должна превышать $24\cdot80000$ часов.

Решая неравенство через сумму геометрической прогрессии получаем, что $n=11$\par

\textbf{Критерии.} \\
\textit{0-4 баллов} $-$ Решение пункта А\\
\textit{0-1 балла} $-$ Составление условия пункта Б \\
\textit{0-2 баллов} $-$ Решение пункта Б\\
\textit{0 баллов} $-$ Отсутствие решения пункта А (сразу НОЛЬ! баллов)\\
\textit{0 баллов} $-$ Любой ответ без объяснений\\

\problem{\textbf{3}} С одной стороны, рыбки каждого цвета должны встречаться по крайней мере в трёх аквариумах из шести, поскольку, если рыбки какого-то цвета лежат не более чем в двух аквариумах, то в оставшихся аквариумах, которых — не менее четырёх, рыбок этого цвета нет. Поэтому всего рыбок должно быть не менее $3\times26$. С другой стороны, если взять по 3 рыбки каждого из 26 цветов и распределить их так, чтобы никакие рыбки одного цвета не были в одном аквариуме, то рыбка любого цвета найдётся в любых четырёх аквариумах (так как остальных аквариумов всего два). Нужное расположение рыбок получится, например, если всех рыбок разбить на 3 группы по 26 рыбок всех цветов и каждую группу распределить поровну (по 13 рыбок) в два аквариума.

\textbf{Критерии.} \\
\textit{7 баллов} $-$ Полностью верное решение\\
\textit{4 баллов} $-$ Верный пример по 13\\
\textit{3 баллов} $-$ Верная оценка, не менее $3\times26$\\
\textit{0-1 балл} $-$ За все остальные попытки решить задачу\\
\textit{0 баллов} $-$ Любой ответ без объяснений\\

\problem{\textbf{4}}
Раскрыв скобки и собрав все члены, содержащие переменные, в левой части неравенства, его нетрудно затем привести к виду

$$
\begin{aligned}
& \left(\frac{a}{b}-1\right)^2+\left(\frac{b}{c}-1\right)^2+\left(\frac{c}{a}-1\right)^2+\left(\frac{a}{b}+\frac{b}{c}+\frac{c}{a}\right)+\left(\frac{a}{c}+\frac{c}{b}+\frac{b}{a}\right) \geq 6 \\
& \frac{a}{b}+\frac{b}{c}+\frac{c}{a} \geq 3 \sqrt[3]{\frac{a}{b} \frac{b}{c} \frac{c}{a}}=3
\end{aligned}
$$



\textbf{Критерии.} \\
\textit{7 баллов} $-$ Решение через неравенство КБШ/Седракяна\\
\textit{7 баллов} $-$ За то что не указано почему сумма квадратов неотрицательна баллы не снижаются\\
\textit{5 баллов} $-$ Верное решение, но не указано почему сумма $\left(\frac{a}{b}+\frac{b}{c}+\frac{c}{a}\right)+\left(\frac{a}{c}+\frac{c}{b}+\frac{b}{a}\right) \geqslant 6$\\
\textit{0-2 балла} $-$  Небольшие продвижения\\ 
\textit{0 баллов} $-$ Любой ответ без объяснений\\

\problem{\textbf{5}}
Отметим точку $F$, симметричную $A$ относительно $E$, и пусть $N$ середина $D F$. Треугольники $B E C$ и $F E D$ подобны по двум сторонам и углу. Отсюда $\angle B M E=\angle F N E$ (как соответствующие элементы) и $\angle F N E=\angle A L E$, поскольку $E N L$ — серединный треугольник $A F D$. Таким образом доказано, что $\angle A L E=\angle B M E$. Прибавим к этому равенству $\angle K L E=\angle K M E$, вычтем из $180^{\circ}$, и получим $\angle K L D=\angle K M C$, откуда $\angle C A D=\angle D B C$, а значит $A B C D$ - вписанный.

\textbf{Критерии.} \\
\textit{7 баллов} $-$ Если решение опирается на расположение точки $F$, баллы не снимаются\\
\textit{3 балла} $-$ Отмечена точка $F$ как в решении\\
\textit{3 баллов} $-$ Доказано подобие треугольников $BEC$ и $FED$\\
\textit{0-1 балл} $-$ За все остальные попытки решить задачу\\

